\documentclass[11pt]{article}
\usepackage[margin=1in]{geometry}
\usepackage{amsmath,amssymb,booktabs,tabularx,longtable}
\usepackage[T1]{fontenc}
\usepackage{lmodern,microtype}
\usepackage[hidelinks]{hyperref}
\hypersetup{pdftitle={Parallax NBA v1.0.0: Pregame Forecasts, Evaluation, and Publication},
  pdfauthor={Eshan Garkhedkar}}
\setlength{\parindent}{0pt}
\setlength{\parskip}{6pt}
\clubpenalty=10000
\widowpenalty=10000
\title{Parallax NBA v1.0.0\\
\large Pregame Forecasts, Evaluation, and Publication}
\author{Eshan Garkhedkar}
\date{}
\begin{document}
\maketitle
\begin{abstract}
Parallax NBA v1.0.0 publishes pregame win probabilities and team ratings
alongside the evidence available at each forecast cutoff. The frozen
forecasting model combines sequential Elo, home court, rest and travel
with a fitted injury correction. Five complete regular seasons provide
6,150 games for the underlying historical archive. The v1.0.0 application
adds a relational SQLite read model and a static Astro interface for
forecasts, results, and evaluation. The forecast model remains v0.3.0.
Public accuracy scores use saved pregame forecasts, while reconstructed
ratings and historical backtests have separate provenance. The contract
calculator gives a model-implied value for a binary payout before fees.
\end{abstract}

\section{The prediction--update cycle}
The model answers two questions in order: how likely is a team to win before
its game, and how should that result change its rating afterward? Ratings
summarize past games. Rest, travel, and home court adjust the next forecast.
Once a game ends, the outcome and eight box-score features determine its
rating update. The new rating is then available for later forecasts.

The sections below follow that order. Each quantity is defined before it is
used in the final update. The eight features $X_1,\ldots,X_8$ are derived in
Section~\ref{sec:features}; they feed a performance factor $P$, which appears
explicitly in the final Elo equation in Section~\ref{sec:final}.
The dictionary in Appendix~\ref{app:symbols} records every symbol and its value
where one is fixed or fitted.

Team $i$ is the home team and $j$ the away team when defining a matchup.
In a team-specific update, the same equations apply to either side using its
own base win probability and outcome. The following is a
step-by-step construction of an empirical model; its adjustments are modeling
choices whose coefficients are fitted, not consequences of a physical law.

\subsection{Methodological basis and scope}
The probability link follows Bradley and Terry's paired-comparison model
\cite{bradleyterry}. With positive strengths $a_i=10^{r_i/c}$ and
$a_j=10^{r_j/c}$, neutral equal-context play gives
$p_i=a_i/(a_i+a_j)$. The complete equation below is the covariate-adjusted
form, with log odds
\begin{equation}
 \ln\!\left(\frac{p_i}{1-p_i}\right)=\frac{\ln 10}{c}(r_i-r_j+Hh)
 -\gamma\ln 10\ln(F_i/G_i).
\end{equation}
Thus the Bradley--Terry probability structure remains, while ratings change
through a sequential Elo-style residual update; the connection between
Elo updates and Bradley--Terry log odds is discussed by Kir\'aly and Qian
\cite{onlineelo}. Home court and the context
term shift log odds before the game.

Whole-History Rating (WHR), described by Coulom \cite{whr}, instead jointly
estimates time-varying rating trajectories under a Bayesian temporal prior
using the available game history. This implementation does not implement
that joint maximum-a-posteriori procedure, backward smoothing, or posterior
rating distributions. Replaying games to fit shared update constants is not
WHR. Its ratings are forward sequential estimates; issued forecasts are
never revised using future outcomes. The schedule, lead and box-score
multipliers below are empirical project specifications, not derivations
from the Bradley--Terry or WHR papers.

\section{Step 0: initialize the season}
The first available historical season starts all teams at 1500. Version 0.3
retains part of the preceding season's final rating deviation:
\begin{equation}\label{eq:carryover}
 r_{\mathrm{start}}=1500+\rho(r_{\mathrm{previous\ final}}-1500),\qquad \rho=0.5.
\end{equation}
A final rating of 1600 therefore becomes 1550 before the new roster is
recentered to a mean of 1500. Season win records,
schedule histories and team performance averages reset. Probability
corrections retain earlier player values. Carryover supplies an initial
strength estimate before the season's roster and results are observed.

\section{Step 1: turn pregame ratings into a win probability}
\subsection{Start with the rating difference and home court}
Let $r_i$ and $r_j$ be the ratings available before tip-off. A basic Elo
forecast compares them on a fixed scale $c$. Home advantage adds $Hh$ to
the home-side rating difference, where $h=1$ at home and $h=0$ at a neutral
venue. With no rest/travel adjustment, the probability would be
\begin{equation}
 p_i=\frac{1}{1+10^{(r_j-r_i-Hh)/c}}.
\end{equation}
A higher home rating or a larger home advantage raises the home win probability.
The away win probability is always $p_j=1-p_i$.

\subsection{Add rest and travel context}
Write $RD$ for rest in intervening calendar days, $TD$ for venue-city travel
in kilometres, and $TZ$ for the time-zone change in hours. Define two positive
context indices:
\begin{align}
 F_i&=RD_j+e^{|TZ_i|/3}+\nu TD_i,\\
 G_i&=RD_i+e^{|TZ_j|/3}+\nu TD_j.
\end{align}
Their ratio compares the two sides' context; it is an empirical adjustment,
not a physiological measurement. The fitted coefficient $\nu$ determines how
much distance contributes. A second coefficient, $\gamma$, determines how
strongly the log ratio affects the forecast. The complete pregame equation is
\begin{equation}\label{eq:probability}
 \boxed{p_i=\frac{1}{1+10^{z_i}}},\qquad
 z_i=\frac{r_j-r_i-Hh}{c}+\gamma\ln(F_i/G_i),\qquad p_j=1-p_i.
\end{equation}
Equal context indices give $\ln(F_i/G_i)=0$, returning the basic Elo forecast.
All inputs to this equation are known before the game. Its outcome and
box score cannot enter their own prediction.

The base is fixed at 10 and $c=497.3558141396255$. Fixing the base and rating scale anchors the model's rating units rather
than fitting redundant scales. Display rounding does not affect calculations.

\subsection{Version 0.3: adjust the published probability for injuries}
Write $p_i$ for the base Elo probability and $q_i$ for the default published
probability. Define $\sigma(a)=(1+e^{-a})^{-1}$ and write the base log odds as $\ln(p/(1-p))$. The correction is an L2-penalized logistic model with the base log odds
as an offset; penalized logistic modeling is discussed in
\cite[Sections 4.4 and 5.5]{esl}. The implemented equation is
\begin{equation}\label{eq:injury}
 \boxed{q_i=\sigma\!\left(\ln\!\left(\frac{p_i}{1-p_i}\right)+b+
 \mathbf w^{\mathsf T}\mathbf d_i\right)},\qquad q_j=1-q_i.
\end{equation}
The vector $\mathbf d_i$ contains standardized pregame covariates: base log
odds, status-weighted prior player minutes and value, and explicit coverage
indicators. It is separate from the postgame features $X_1,\ldots,X_8$ below.
Means, standard deviations, intercept and weights use earlier seasons only
and are frozen with a feature-order schema.

Official Out, Doubtful, Questionable and Probable designations are separate
\cite{injuries}. They are categorical predictors, not fixed probabilities
of playing or clinician estimates.
For each player, estimate expected minutes $\bar m$ with a normalized
exponential mean over earlier played observations, half-life ten observations.
The descriptive player score uses the Game Score formula documented by
Sports Reference \cite{gamescore}:
\begin{align}
 \mathrm{PlayerScore}={}&PTS+0.4FGM-0.7FGA-0.4(FTA-FTM)\notag\\
 &+0.7OREB+0.3DREB+STL+0.7AST+0.7BLK-0.4PF-TOV.
\end{align}
$BLK$ and $PF$ are blocks and personal fouls. With discounted prior sums
$V$ of PlayerScore and $M$ of minutes, the shrunk value per minute is
\begin{equation}
 v_p=\frac{V+80(0.3)}{M+80}.
\end{equation}
For team $t$ and status $s$, define the two contributions
\begin{equation}
 A_{t,s}=\sum_{a\in(t,s)}\frac{\bar m_a}{240},\qquad
 B_{t,s}=\sum_{a\in(t,s)}\frac{\bar m_a v_a}{80}.
\end{equation}
The eight status features are $A_{i,s}-A_{j,s}$ and $B_{i,s}-B_{j,s}$
for the four statuses. Ten further coverage features comprise differences
in submitted flags, unmapped counts, missing-history counts and valued-player
counts, plus each side's submitted flag, unmapped count and missing-history
count separately. Together with base log odds, the correction has 19 inputs.
For raw input $u_k$, the fitted standardization is
$d_k=(u_k-\mu_k)/s_k$, where $\mu_k,s_k$ use training rows only and
scales below $10^{-10}$ are replaced by one. The base log odds appear both as a fixed
offset and as a standardized covariate, allowing fitted recalibration. The constants 240 and
80 set feature scales; 80 minutes at value 0.3 supplies a player-value prior.
Coverage features distinguish submitted, unsubmitted, unmapped and
no-prior-history cases. Missing records never imply observed health or zero
playing time. These descriptive player values are not RAPM.

If no verified report covers the game, use $q_i=p_i$. Partially submitted
reports use explicit unknown indicators. Live collection rejects future-dated
or more-than-16-hour-old reports. The Elo update always uses $p_i$, not $q_i$:
the injury correction does not feed back into ratings.

\subsubsection{From source reports to issued forecasts}
The scheduled collector checks the newest eligible PDF listed on the
official index, when available. It can also check up to four already-due
public quarter-hour versions, skipping duplicate URLs. It parses the PDF table and actual header publication time,
rejects future or more-than-16-hour-old documents, and matches game date,
teams and matchup. Normalized player names map only when unambiguous to
ESPN athlete identifiers. Earlier completed ESPN player box scores supply
minutes and PlayerScore; invalid team/player totals are excluded from player
history rather than interpreted as zero minutes \cite{espn,sportsdataverse}.
Expected-minute smoothing, shrinkage constants and feature scales are project choices, not independently established player-impact
estimates or injury severity measurements.

Historical forecasts freeze the whole Eastern-calendar slate at the earlier
of 10 a.m. and one hour before its first tip-off; no player performance from
that slate enters its own morning forecasts. Live revisions use information
collected at their actual issuance time, never a backdated timestamp.
Each issued revision preserves the model fingerprint and coverage. When
a verified report is used, it also records the report URL, publication
time, hash, and numerical injury inputs. Partial submissions
use unknown indicators; absent verified matchup coverage gives $q_i=p_i$.
The two-hour collection cadence can lag later NBA report changes, which is
why the dashboard exposes the report timestamp. Game results update the Elo
state through the base residual $y_i-p_i$; the published injury probability
$q_i$ is used for forecasts and their evaluation, not that rating residual.

\section{Step 2: derive eight features from the completed box score}\label{sec:features}
Only after the game finishes do its box-score counts become inputs to the
rating update. Let $O$ denote the opponent. First estimate possessions:
\begin{equation}\label{eq:possessions}
 \mathrm{Poss}=FGA+0.44FTA-OREB+TOV.
\end{equation}
$FGA$ and $FTA$ are field-goal and free-throw attempts, $OREB$ is offensive
rebounds, and $TOV$ is summed player turnovers. The coefficient 0.44
approximates the possession cost of free throws; offensive rebounds are
subtracted because they extend a possession. Opponent possessions use the
same expression with opponent counts.

The eight features describe scoring efficiency, shooting, rebounding,
ball movement, ball security, and steals:
\begin{align}
 X_1&=\frac{100(PTS-PTS_O)}{\mathrm{Poss}}
 &&\text{net efficiency},\\
 X_2&=\frac{FGM+0.5(3PM)}{FGA}
 &&\text{effective field-goal percentage},\\
 X_3&=\frac{PTS}{2(FGA+0.44FTA)}
 &&\text{true shooting percentage},\\
 X_4&=\frac{OREB}{OREB+DREB_O}
 &&\text{offensive-rebound rate},\\
 X_5&=\frac{DREB}{DREB+OREB_O}
 &&\text{defensive-rebound rate},\\
 X_6&=\frac{AST}{\max(TOV,1)}
 &&\text{assist-to-turnover ratio},\\
 X_7&=\frac{TOV}{FGA+0.44FTA+TOV}
 &&\text{turnover fraction},\\
 X_8&=\frac{STL}{\mathrm{Poss}_O}
 &&\text{steal rate}.
\end{align}
Here $PTS$ means points, $FGM$ means field goals made, $3PM$ means three-point
makes, $DREB$ means defensive rebounds, $AST$ means assists, and $STL$ means
steals. Rates are fractions rather than percentages on a 0--100 scale.
$X_1$ measures the scoring margin per 100 estimated team possessions. The two rebound
rates pair each team category with the opposite opponent category. A
zero-turnover game uses one in the assist-ratio denominator; other
nonpositive denominators fail validation.

\section{Step 3: combine the features into a performance factor}\label{sec:performance}
\subsection{Scale and weight the eight features}
The raw features have different units and ranges. Denote their fixed scaled
versions by lowercase $x$:
\begin{equation}\label{eq:scaled-features}
 (x_1,\ldots,x_8)=(X_1/10,X_2,X_3,X_4,X_5,X_6/3,-X_7,X_8).
\end{equation}
Dividing net efficiency by 10 and assist ratio by 3 sets their relative scale.
The negative sign on turnover fraction makes fewer turnovers contribute a
larger score. These scales were fixed before fitting.

The fitted weights $\alpha_1,\ldots,\alpha_8$ are positive and sum to one.
The weighted box-score total is
\begin{equation}\label{eq:game-score}
 \mathrm{Score}=\sum_{k=1}^{8}\alpha_k x_k.
\end{equation}
Thus every $X_k$ affects the total through its scaled value $x_k$ and its
weight $\alpha_k$. A large weight is not itself a causal claim or a scale-free
measure of feature importance.

\subsection{Compare this game with the team's previous games}
Convert the weighted total into a bounded game-performance score:
\begin{equation}\label{eq:game-perf}
 \mathrm{Perf}=\operatorname{clip}
 \left(\frac{1}{1+e^{-\mathrm{Score}}},10^{-6},1-10^{-6}\right).
\end{equation}
The sigmoid avoids singularities and accepts negative efficiency values.
Clipping limits its result to the stated lower and upper bounds.
For a team with $n$ earlier games, its reference performance is
\begin{equation}\label{eq:prior-perf}
 \mathrm{PriorPerf}=\frac{1}{n}\sum_{g=1}^{n}\mathrm{Perf}_g.
\end{equation}
This average excludes the current game. The relative-performance factor is
\begin{equation}\label{eq:performance-factor}
 \boxed{P=\operatorname{clip}
 \left(\frac{\mathrm{Perf}}{\mathrm{PriorPerf}},0.25,4\right)}.
\end{equation}
$P>1$ means better box-score performance than the team's prior average;
$P<1$ means worse. Season openers use $P=1$, because no previous team games
are available. Stored data calls $\mathrm{Perf}$ the aggregate game performance
factor (AGPF) and $P$ the performance adjustment factor (PAF).

The complete feature path is therefore
\begin{equation}\label{eq:feature-path}
 X_1,\ldots,X_8\ \longrightarrow\ x_1,\ldots,x_8
 \ \longrightarrow\ \mathrm{Score}
 \ \longrightarrow\ \mathrm{Perf}\ \longrightarrow\ P.
\end{equation}
This factor is now ready to enter the Elo update; it is not an additional
input to the already-issued probability for the same game.

\section{Step 4: assemble the final Elo update}\label{sec:final}
\subsection{Start with the outcome relative to the prediction}
For the team being updated, let $y=1$ for a win and $y=0$ for a loss, and let
$p$ be its pregame probability from Equation~\eqref{eq:probability}. The
basic Elo update is
\begin{equation}
 r_{\mathrm{new}}=r+K(y-p).
\end{equation}
$K$ controls update size. An unexpected win produces a larger positive error
than an expected win; a loss gives a negative error. The model multiplies this
base update by three adjustments. Define the result sign $s=2y-1$, equal to
$+1$ for a win and $-1$ for a loss.

\subsection{Adjustment 1: strength of schedule}
A team's prior smoothed win rate is
\begin{equation}
 W_t=\frac{w_t+1}{n_t+2},
\end{equation}
where $w_t$ and $n_t$ count only wins and games before the forecast cutoff.
Let $\overline W_{\mathrm{opp\ history}}$ be the mean smoothed win rate of
its opponent's previous opponents, with repeated opponents counted in
proportion to their frequency. Then
\begin{equation}
 S=\frac{2W_{\mathrm{opp}}+\overline W_{\mathrm{opp\ history}}}{3},
 \qquad M_S=1+s\beta S.
\end{equation}
Empty histories use 0.5. The fitted coefficient $\beta$ controls this
schedule multiplier. Its value is calculated from prior results, not the
current game's box score.

\subsection{Adjustment 2: lead retention}
Let $LL$ and $LL_O$ be the two teams' largest leads, and $PTS$ and $PTS_O$
their final scores. Define
\begin{equation}
 L=\frac{|PTS-PTS_O|}{\max(LL,LL_O)},
 \qquad M_L=1+\phi\log_2(1+L).
\end{equation}
For valid non-tied games, $0<L\leq1$. This compares the final margin with the
largest advantage held by either side. It remains defined even when one team
never led. Peak leads below the final margin fail validation. The fitted
coefficient $\phi$ sets the influence of this adjustment.

\subsection{Adjustment 3: box-score performance}
Use the factor $P$ already derived from the eight features:
\begin{equation}
 M_P=1+s\lambda(P-1).
\end{equation}
For a winner, $P>1$ increases the positive update and $P<1$ reduces it. For
a loser, $P>1$ reduces the magnitude of its negative update and $P<1$
increases it. With $P$ clipped to $[0.25,4]$ and
$0\leq\lambda\leq0.25$, the multiplier cannot become negative.
This is exactly where the $X$ features enter Elo.

\subsection{The complete equation}
Multiplying the three adjustments gives the final model:
\begin{equation}\label{eq:update}
 \boxed{r_{\mathrm{new}}=r+K(y-p)\,M_S\,M_L\,M_P}.
\end{equation}
Substituting their definitions makes every contribution explicit:
\begin{equation}\label{eq:update-expanded}
 \boxed{r_{\mathrm{new}}=r+K(y-p)
 (1+s\beta S)(1+\phi\log_2(1+L))(1+s\lambda(P-1))}.
\end{equation}
The probability $p$ uses prior information; $L$ and $P$ use the just-completed
game. The dependence of the final multiplier on the box score is
\begin{equation}
 (X_1,\ldots,X_8)\longrightarrow P\longrightarrow M_P
 \longrightarrow r_{\mathrm{new}}\longrightarrow
 \text{later-game probabilities}.
\end{equation}
Both teams' updates use their pregame ratings and are computed before either
is committed. Asymmetric updates can change the league mean. After the
historical slate, subtracting that mean minus 1500 from every rating recenters
the league without changing pairwise rating differences. Live state recenters
after individual completed games. Inter-season initialization follows
Equation~\eqref{eq:carryover}; it does not change the within-season equation.

\section{Data construction}
Free ESPN-derived schedule, team-box-score, and play-by-play parquet files from
SportsDataverse provide the initial history \cite{sportsdataverse}. Primary
ESPN game summaries replace incomplete or inconsistent archival records
\cite{espn}. GeoNames supplies venue-city coordinates and IANA time zones
\cite{geonames}. Each of the five seasons contains 1,230 completed standard
regular-season games, 30 teams, and 82 games per team. Preseason, playoffs,
All-Star games, and the NBA Cup final are excluded; Cup games that count
toward the regular season remain included.

Validation requires unique game/team identities, matching opponent rows and
header scores, finite nonnegative integer counts, made shots no greater than
attempts, scoring identities, plausible assists and steals, valid feature
denominators, and peak leads consistent with final margins. Missing values
are not converted to zero. Largest leads can be recovered from a score path
only when its final scores match the header. Inconsistent team aggregates
can be replaced by complete player rows with matching points and shooting
identities. Turnovers consistently mean summed player turnovers, excluding
provider non-player team turnovers.

\begin{table}[ht]
\centering\small
\begin{tabular}{lrrr}
\toprule
Season & Accepted games & Primary replacements & Unresolved failures\\
\midrule
2021--22 & 1230 & 132 & 0\\
2022--23 & 1230 & 158 & 0\\
2023--24 & 1230 & 201 & 0\\
2024--25 & 1230 & 165 & 0\\
2025--26 & 1230 & 217 & 0\\
\bottomrule
\end{tabular}
\caption{Source audit. The 873 primary replacements mostly address incomplete
archived play-by-play; they are not a count of corrupt final box scores.}
\end{table}

All required numerical formula inputs are complete. Repair actions, availability
proxies, source URLs, retrieval times where available, and SHA-256 hashes are
stored with the data or collection manifests. This establishes internal
consistency, not perfect measurement of every real-world quantity. Travel is
estimated between successive venue cities, not actual itineraries. Rest counts
intervening Eastern dates, so back-to-backs have zero. Openers receive three
rest days and zero travel/time-zone change, with an initialization flag.
Both venue offsets are evaluated at the current tip-off instant, preventing
a daylight-saving change alone from creating geographic travel. Ambiguous
or unknown cities cause validation failure.

\subsection{Chronology and forecast availability}
A prediction uses only information available at its recorded forecast cutoff.
The home-win outcome is observed afterward.
The historical replay freezes all forecasts for an Eastern-calendar daily
slate at the earlier of 10:00 America/New\_York and one hour before its first
tip-off. It requires all prior slates to have completed before that cutoff;
current-slate results cannot enter its forecasts. Ratings update after the
slate. A team cannot appear twice on one slate, and duplicate games are rejected.

Completion availability is proxied by terminal EndGame wallclock timestamps,
rather than later edited play timestamps. Seventeen games use dated final-recap
publication times. One game with source play dates shifted by a day uses an
explicitly flagged conservative tip-off-plus-eight-hours availability bound;
this is not an observed completion time. Replays check these proxies against
the next cutoff but cannot establish when historical corrections first reached
a live feed. Genuinely issued forecasts record their actual issue time,
data cutoff, and model version in a separately preserved ledger.


\section{Frozen model fitting and evaluation}
The v0.3.0 core fit uses five complete regular seasons. Teams enter the
fitting objective after both have ten prior games. Earlier games still
update ratings and histories. SciPy differential evolution with L-BFGS-B
polishing minimizes mean binary log loss over the 5,370 eligible games
\cite{scipy}. The fit converged in 37 generations at loss 0.615970.
Seven free feature-weight coordinates and a fixed eighth coordinate
produce eight normalized weights via softmax. The lead and performance
coefficients reached their fitting bounds.

For the injury correction, historical base probabilities are generated
with earlier-prefix fits, with fixed initial defaults for the first
anchor season. The injury head minimizes mean binary log loss plus
$\eta\|\mathbf w_*\|_2^2/2$, with the intercept included in
$\mathbf w_*$. Earlier validation selects the penalty from
$\{10^{-4},10^{-3},10^{-2},0.1,1,10\}$ and an unchanged-base option.
The frozen head uses $\eta=10^{-4}$. Standardization is learned from
fit rows only. Full coefficients and scales are stored in the frozen bundle.

Across 816 historical game dates, 815 official morning PDFs were retrieved
and validated. Player totals and minutes pass checks for 12,235 of
12,300 team-games. Invalid player histories are excluded while valid team
records remain usable. Reports must match date, teams and matchup and
precede the saved cutoff.

\section{Reading the accuracy results}
The published Accuracy page scores one compatible saved pregame forecast
per completed game. The matchup cutoff filter includes only games whose
recorded completion time precedes the selected forecast's saved model
cutoff. Users can filter by team, season, and whether both teams had ten
prior games. Training replays measure historical fit on development seasons.
Post-fit replays apply the same frozen bundle to later games using saved
inputs. The analytics database keeps those rows and chronological backtests
separate from original issued forecasts.

\begin{table}[ht]
\centering\small
\begin{tabular}{lrrr}
\toprule
Procedure & Log loss & Brier score & Winner accuracy\\
\midrule
Carryover Elo & 0.603496 & 0.208434 & 66.64\%\\
Elo with injury correction & 0.598309 & 0.206217 & 67.07\%\\
Boosted comparison & 0.600005 & 0.206871 & 67.15\%\\
\bottomrule
\end{tabular}
\caption{Matched retrospective comparisons on 3,690 games from
2023--24 through 2025--26, including season openers.}
\end{table}

The matched comparison \cite{evaluation} fits and selects on earlier
seasons, then evaluates the next season. These outcomes were already
studied during development, so the scores are retrospective evidence.
The boosted comparator is recorded for reproducibility. The public
forecast is the Elo model with the injury correction and its missing-report
fallback.

Overall calibration groups games by the predicted probability of the
favorite, choosing home at exactly $0.5$. Five-percentage-point bins
include the lower boundary and exclude the upper one, except the final
bin includes one. Each bin shows the game count, mean forecast, observed
win rate and Wilson 95\% interval \cite{wilson}. Team-specific
calibration instead uses the selected team's win probability in
ten-percentage-point bins. Expected wins for a team are
$\sum_j q_j$ over its selected appearances. Winner accuracy counts
the higher-probability team winning. Brier score and log loss evaluate
the assigned probabilities; both scores are lower when better.

\section{Forecast records and reconstruction}
Each original issued revision keeps its observed issuance time, recorded
tip, probabilities, model and parameter fingerprints, and captured
inputs. Changed pregame inputs append a revision. The live scorecard
selects the latest matching forecast issued before the valid game start.
Backfilled outcomes update the rating state and remain outside the
issued scorecard when an original forecast is unavailable.

Current-method replays apply the frozen model to stored pregame inputs.
Historical replays freeze the daily Eastern slate before updating the
same slate's ratings and player state. Current-season reconstructions
use the saved pregame cutoff and captured schedule context where
available. A replay can reuse captured numerical injury covariates for
an older revision. Without verified inputs the injury head falls back to
base Elo. The historical records incorporate provider corrections, so a
reconstruction has a source-vintage limitation even with a pregame
cutoff. Bundle and source-result fingerprints identify compatible
replay results. Matchup pages show the compatible saved forecast where
available and explicitly flag games with no saved pregame forecast.
Reconstructed ratings support the team histories. The public matchup and
accuracy records use the original issued ledger.

\section{V1 data and publication architecture}
The Python collector validates NBA schedule and result records and saves
the season archive. The archive remains the authoritative record of
forecast revisions and outcomes. The analytics build imports one
coherent archive into SQLite. Games, team-game observations, forecast
revisions, rating histories and evaluation runs have explicit keys.
A window query selects the latest compatible original forecast per
completed game for scoring. An export allowlist limits the public site to
approved fields, including team-level injury coverage and saved forecasts.
The build checks source and export fingerprints. The database and
raw provider caches stay outside the website artifact.

Astro builds the website from the validated exports and serves dated
slates, matchup pages, an archive, team rating histories, accuracy
filters, and model documentation. GitHub Actions is scheduled every two
hours and can also run on demand. It saves the updated ledger, builds
from that exact commit, and runs browser checks before publishing to
\href{https://eshangark.github.io/parallax-nba/}{GitHub Pages}. If the
pipeline fails, the last successful site remains available. The public
\href{https://eshangark.github.io/parallax-nba/release.json}{release
metadata} records the application and model versions, source commit,
and last successful collection timestamp. A forecast's issuance time
records collection and model execution. Public availability can lag
that time while the site is built and deployed.

\section{Contract pricing before fees}
For a contract paying one dollar if the selected team wins and zero
otherwise, let $q$ denote the selected team's forecast and let $a$
denote the entered purchase price in cents. Under normal binary
settlement the model-implied value, expected profit and break-even
probability are
\begin{align}
 V_{\mathrm{model}} &= 100q, &
 \operatorname{E}[\mathrm{profit}] &= 100q-a, &
 q_{\mathrm{break-even}} &= a/100.
\end{align}
The realized profit in cents is $100-a$ on a win and $-a$ on a loss.
For example, $q=0.70$ and $a=66$ give an expected profit of four cents
per contract under the model. The calculator assumes the entered price
is the execution price, excludes fees, and evaluates holding through
settlement. A market comparison would require timestamped executable
quotes and the applicable market terms and costs. The predictive
benchmarks measure forecast quality.

\section{Reproducibility}
The repository includes source fingerprints, validation records,
frozen bundles and the matched-sample audit. The historical audit
checks exact game identity, cutoff, outcome and eligibility before
comparing models. Numerical parity checks allow differences within
$10^{-12}$ while preserving forecast identity and record structure.
The local analytics build can reproduce the site data from a pinned
archive without fetching additional provider data. The repository
documents the build and refresh commands.

\appendix
\section{Travel and time-zone derivation}
For venue latitude $\varphi$ and longitude $\ell$ in radians,
\begin{align}
 a_v&=\sin^2((\varphi_2-\varphi_1)/2)
 +\cos\varphi_1\cos\varphi_2\sin^2((\ell_2-\ell_1)/2),\\
 TD&=2R\arcsin\sqrt{\operatorname{clip}(a_v,0,1)},\qquad R=6371\ \text{km}.
\end{align}
Neutral-game venue cities are used rather than substituting a franchise home
city. If $o(v,t)$ is a venue's dated UTC offset in hours, then
$TZ=o(v_{\mathrm{current}},t_g)-o(v_{\mathrm{previous}},t_g)$.
Exact arenas and actual itineraries are not measured by these proxies.

\clearpage
\section{Symbol dictionary}\label{app:symbols}
The dictionary lists model constants and the quantities used in the construction. Descriptive
names replace opaque intermediate symbols: Score is the weighted feature
sum, Perf is a bounded game score, PriorPerf is its prior average, and P is
their clipped ratio. Constants are rounded for display; the frozen release
configuration retains full precision. Dictionary values refer to the current
version 0.3.0 core fit.
\begingroup
\normalsize
\setlength{\tabcolsep}{4pt}
\renewcommand{\arraystretch}{1.14}
\subsection*{Fitted constants and fixed scales}
\begin{longtable}{@{}>{\raggedright\arraybackslash}p{0.23\textwidth}>{\raggedright\arraybackslash}p{0.51\textwidth}>{\raggedright\arraybackslash}p{0.22\textwidth}@{}}
\toprule
Symbol & Meaning & Value or range\\
\midrule
\endfirsthead
\toprule
Symbol & Meaning & Value or range\\
\midrule
\endhead
\bottomrule
\endfoot
$K$ & The base update coefficient controls how much a game can change Elo. & 13.748881\\
$H$ & Home advantage adds this many Elo points to the home-side rating difference. & 40.03439907\\
$c$ & The fixed rating scale converts Elo differences into probability differences. & 497.3558141 (fixed)\\
$\gamma$ & The fatigue coefficient weights the logarithm of the rest/travel index ratio. & 0.179559595\\
$\nu$ & The travel coefficient weights venue-city distance inside the fatigue indices. & 9.534356884e-05\\
$\beta$ & The schedule coefficient weights the prior strength-of-schedule adjustment. & 0.7081901219\\
$\phi$ & The lead coefficient weights the logarithmic lead-retention adjustment. & 2\\
$\lambda$ & The performance coefficient weights relative box-score performance. & 0.25\\
$10$ & The probability exponent uses a fixed base of ten. & 10 (fixed)\\
$3$ & The fatigue indices divide absolute time-zone change by this fixed hourly scale. & 3 hours (fixed)\\
$R$ & The Haversine calculation uses this fixed Earth-radius approximation. & 6371 km\\
\end{longtable}
\subsection*{Learned feature weights}
\begin{longtable}{@{}>{\raggedright\arraybackslash}p{0.23\textwidth}>{\raggedright\arraybackslash}p{0.51\textwidth}>{\raggedright\arraybackslash}p{0.22\textwidth}@{}}
\toprule
Symbol & Meaning & Value or range\\
\midrule
\endfirsthead
\toprule
Symbol & Meaning & Value or range\\
\midrule
\endhead
\bottomrule
\endfoot
$\alpha_1$ & The normalized weight multiplies scaled net efficiency. & 0.0006686328527\\
$\alpha_2$ & The normalized weight multiplies scaled effective field-goal percentage. & 0.0008033843254\\
$\alpha_3$ & The normalized weight multiplies scaled true shooting percentage. & 0.008884867219\\
$\alpha_4$ & The normalized weight multiplies scaled offensive-rebound rate. & 0.0007283531314\\
$\alpha_5$ & The normalized weight multiplies scaled defensive-rebound rate. & 0.1613285001\\
$\alpha_6$ & The normalized weight multiplies scaled assist-to-turnover ratio. & 0.8047679813\\
$\alpha_7$ & The normalized weight multiplies scaled turnover fraction. & 0.01518263624\\
$\alpha_8$ & The normalized weight multiplies scaled steal rate. & 0.007635644786\\
\end{longtable}
\subsection*{Matchup and rating-update quantities}
\begin{longtable}{@{}>{\raggedright\arraybackslash}p{0.23\textwidth}>{\raggedright\arraybackslash}p{0.51\textwidth}>{\raggedright\arraybackslash}p{0.22\textwidth}@{}}
\toprule
Symbol & Meaning & Value or range\\
\midrule
\endfirsthead
\toprule
Symbol & Meaning & Value or range\\
\midrule
\endhead
\bottomrule
\endfoot
$i,j,O$ & In the probability equation, i is the home team and j the away team; O denotes an opponent in feature derivations. & Team identifiers\\
$g,t$ & The index g labels a game; t labels a team in win-rate notation and a time instant in the offset function. & Context dependent\\
$r_i,r_j$ & These are the team ratings available before the game. & Carryover; first anchor 1500\\
$r_{\mathrm{new}}$ & This is a team rating after its game update, before league recentering. & Varies by game\\
$p_i,p_j$ & These are the complementary pregame probabilities of a home and away win. & Between 0 and 1; sum 1\\
$z_i$ & The probability exponent combines the rating difference, home advantage, and fatigue adjustment. & Varies by matchup\\
$h$ & The venue indicator enables home advantage at a home venue and removes it at a neutral venue. & 1 home; 0 neutral\\
$RD_i,RD_j$ & Rest is the number of intervening Eastern-calendar dates since the previous game. & Varies; opener 3 days\\
$TD_i,TD_j$ & Travel is the estimated distance between consecutive venue cities. & Varies; opener 0 km\\
$TZ_i,TZ_j$ & Time-zone change compares dated venue offsets at the current tip-off. & Varies; opener 0 hours\\
$F_i,G_i$ & The two empirical indices combine each side's rest, travel, and time-zone context. & Varies; positive\\
$y_i\;(y)$ & The observed outcome is one for a win and zero for a loss. & 1 win; 0 loss\\
$s_i\;(s)$ & The result sign equals twice the win indicator minus one. & +1 win; -1 loss\\
$S_i\;(S)$ & Schedule strength combines the opponent's prior win rate and its prior opponents' win rates. & Varies; empty history 0.5\\
$W_t$ & The smoothed prior win rate equals (prior wins + 1) divided by (prior games + 2). & Varies; no games 0.5\\
$w_t,n_t$ & These are the wins and games observed for a team before the forecast cutoff. & Nonnegative integer counts\\
$\overline W_{\mathrm{opp\ history}}$ & This is the mean prior win rate of the opponent's previous opponents, counting repeats. & Varies; empty mean 0.5\\
$LL_i,LL_j$ & Largest lead is the maximum points advantage achieved by each team during the completed game. & Nonnegative point counts\\
$L_i\;(L)$ & Lead retention divides absolute final margin by the larger of both teams' largest leads. & Greater than 0; at most 1\\
$\mathrm{Score}$ & The box-score total is the weighted sum of the scaled box-score features. & Varies by team-game\\
$\mathrm{Perf}$ & The aggregate game performance factor (AGPF) applies a clipped sigmoid to the weighted Score. & From 0.000001 to 0.999999\\
$\mathrm{PriorPerf}$ & The baseline is the expanding mean of the team's prior game-performance values. & Varies; opener uses current Perf\\
$P$ & The performance adjustment factor (PAF) compares current game performance with its prior baseline and clips the ratio. & From 0.25 to 4; opener 1\\
$n$ & The prior-performance average divides the sum of a team's earlier game-performance scores by this count. & Number of prior team games\\
$M_S$ & The schedule multiplier equals one plus the result sign times beta times schedule strength. & Varies by game\\
$M_L$ & The lead multiplier equals one plus phi times the base-two logarithm of one plus lead retention. & Varies by game\\
$M_P$ & The performance multiplier equals one plus the result sign times lambda times (P minus one). & Varies by game\\
$\Delta r$ & The rating change equals the base outcome error multiplied by the three adjustment factors. & Varies by game\\
\end{longtable}
\clearpage
\subsection*{Box-score features and counts}
\begin{longtable}{@{}>{\raggedright\arraybackslash}p{0.23\textwidth}>{\raggedright\arraybackslash}p{0.51\textwidth}>{\raggedright\arraybackslash}p{0.22\textwidth}@{}}
\toprule
Symbol & Meaning & Value or range\\
\midrule
\endfirsthead
\toprule
Symbol & Meaning & Value or range\\
\midrule
\endhead
\bottomrule
\endfoot
$X_1$ & The feature measures point differential per 100 estimated team possessions. & Varies by team-game\\
$X_2$ & The feature credits three-point makes with half an additional made shot. & Varies by team-game\\
$X_3$ & The feature measures points per twice the estimated shooting opportunities. & Varies by team-game\\
$X_4$ & The feature divides offensive rebounds by offensive rebounds plus opponent defensive rebounds. & Varies by team-game\\
$X_5$ & The feature divides defensive rebounds by defensive rebounds plus opponent offensive rebounds. & Varies by team-game\\
$X_6$ & The feature divides assists by player turnovers, using a denominator of one for zero turnovers. & Varies by team-game\\
$X_7$ & The feature divides player turnovers by field-goal attempts plus 0.44 free-throw attempts plus player turnovers. & Varies by team-game\\
$X_8$ & The feature divides steals by estimated opponent possessions. & Varies by team-game\\
$x_k$ & The scaled features follow (X1/10, X2, X3, X4, X5, X6/3, -X7, X8). & Fixed transformation\\
$\mathrm{Poss},\mathrm{Poss}_O$ & Estimated possessions equal field-goal attempts plus 0.44 free-throw attempts minus offensive rebounds plus player turnovers. & Varies; coefficient 0.44\\
$PTS,PTS_O$ & These counts record points scored by the team and its opponent. & Points in the completed game\\
$FGM,FGA$ & These counts record field goals made and attempted. & Nonnegative integer counts\\
$3PM$ & This count records made three-point field goals. & Nonnegative integer count\\
$FTA$ & This count records free throws attempted. & Nonnegative integer count\\
$OREB,DREB$ & These counts record offensive and defensive rebounds; an O subscript denotes opponent counts. & Nonnegative integer counts\\
$AST,TOV,STL$ & These counts record assists, summed player turnovers, and steals. & Nonnegative integer counts\\
\end{longtable}
\clearpage
\subsection*{Fitting, evaluation, and travel notation}
\begin{longtable}{@{}>{\raggedright\arraybackslash}p{0.23\textwidth}>{\raggedright\arraybackslash}p{0.51\textwidth}>{\raggedright\arraybackslash}p{0.22\textwidth}@{}}
\toprule
Symbol & Meaning & Value or range\\
\midrule
\endfirsthead
\toprule
Symbol & Meaning & Value or range\\
\midrule
\endhead
\bottomrule
\endfoot
$\theta$ & The fitted parameter vector contains seven scalar coefficients and seven independent feature-weight logits. & 14 free coordinates\\
$\mathcal L(\theta)$ & Training loss is the mean binary log loss across eligible training games. & 0.6159695469 at fitted solution\\
$N$ & The number of eligible games normalizes a stated fitting or evaluation score. & 5370 core fit\\
$\varphi_1,\varphi_2,\ell_1,\ell_2$ & These are the latitudes and longitudes of consecutive venue cities, expressed in radians. & Varies by venue\\
$a_v$ & The Haversine intermediate combines the angular latitude and longitude differences. & Between 0 and 1\\
$v,o(v,t)$ & The venue v and its UTC-offset function define a location's dated time-zone offset. & Venue dependent; hours\\
$\ln,\log_2,\exp$ & These functions are the natural logarithm, base-two logarithm, and natural exponential. & Function definitions\\
$\operatorname{clip}$ & Clipping limits an input to its stated lower and upper bounds. & Bounds stated above\\
\end{longtable}

\clearpage
\subsection*{Carryover and probability corrections}
\begin{longtable}{@{}>{\raggedright\arraybackslash}p{0.23\textwidth}>{\raggedright\arraybackslash}p{0.51\textwidth}>{\raggedright\arraybackslash}p{0.22\textwidth}@{}}
\toprule
Symbol & Meaning & Value or range\\
\midrule
\endfirsthead
\toprule Symbol & Meaning & Value or range\\\midrule
\endhead
\bottomrule\endfoot
$\rho$ & The season-start rating retains this fraction of the prior final deviation from 1500. & 0.5\\
$q_i,q_j$ & These are the default published injury-adjusted probabilities; their sum is one. & Between 0 and 1\\
$\sigma,\ln(p/(1-p))$ & The sigmoid and log odds convert between probability and an unconstrained score. & Defined in the text\\
$\mathbf d_i$ & The standardized injury-head vector contains base log odds and prior player/report features. & 19 coordinates\\
$b,\mathbf w$ & The intercept and standardized-feature weights define the fitted logistic correction. & b = 0.00231409; weights in JSON\\
$\mathbf w_*,\eta$ & The penalized vector includes the intercept; eta controls its squared-norm penalty. & 20 coefficients; eta 0.0001\\
$\bar m$ & Expected minutes average a player's earlier played observations with exponential decay. & Half-life 10 observations\\
$a_i,a_j$ & Positive paired-comparison strengths are powers of ten of the rating divided by its scale. & $10^{r_i/c},10^{r_j/c}$\\
$A_{t,s},B_{t,s}$ & Team $t$'s status $s$ totals summarize expected minutes and expected minutes times player value. & Vary by report\\
$u_k,\mu_k,s_k,d_k$ & A raw injury-head input, its training mean and scale, and its standardized value determine one correction term. & $d_k=(u_k-\mu_k)/s_k$\\
$\mathrm{PlayerScore}$ & The descriptive player box-score total combines scoring, shooting, rebounds, assists, steals, blocks, fouls and turnovers. & Formula in the text\\
$M,V$ & Discounted sums retain earlier player minutes and PlayerScore. & Half-life 10 observations\\
$v_p$ & Prior value per minute is shrunk toward a fixed descriptive reference. & 80 minutes at value 0.3\\
$FTM,BLK,PF$ & These counts are free throws made, blocks and personal fouls. & Completed player box only\\
\end{longtable}

\endgroup

\clearpage
\begingroup
\raggedright
\begin{thebibliography}{99}

\bibitem{evaluation}
Parallax NBA. Matched historical evaluation and reproducibility notes.
\url{https://github.com/eshangark/parallax-nba/blob/main/docs/EVALUATION.md}.
\bibitem{bradleyterry}
R. A. Bradley and M. E. Terry (1952).
Rank analysis of incomplete block designs: I. The method of paired comparisons.
\emph{Biometrika}, 39(3--4), 324--345.
\href{https://doi.org/10.1093/biomet/39.3-4.324}{doi:10.1093/biomet/39.3-4.324}.
\bibitem{whr}
R. Coulom (2008). Whole-History Rating: A Bayesian rating system for players
of time-varying strength. In \emph{Computers and Games}, LNCS 5131,
113--124. Springer. \url{https://www.remi-coulom.fr/WHR/}.
\bibitem{onlineelo}
F. J. Kir\'aly and Z. Qian (2017). Modelling competitive sports:
Bradley--Terry--Elo models for supervised and on-line learning of paired
competition outcomes. \emph{arXiv:1701.08055}.
\url{https://arxiv.org/abs/1701.08055}.
\bibitem{gamescore}
Sports Reference. \emph{Basketball-Reference glossary}, Game Score definition.
\url{https://www.basketball-reference.com/about/glossary.html}.
\bibitem{esl}
T. Hastie, R. Tibshirani and J. Friedman (2009).
\emph{The Elements of Statistical Learning}, second edition. Springer.
Sections 4.4 and 5.5: logistic regression and regularization.
\url{https://hastie.su.domains/ElemStatLearn/}.
\bibitem{wilson}
E. B. Wilson (1927). Probable inference, the law of succession, and statistical inference.
\emph{Journal of the American Statistical Association}, 22(158), 209--212.
\href{https://doi.org/10.1080/01621459.1927.10502953}{doi:10.1080/01621459.1927.10502953}.
\bibitem{scoring}
T. Gneiting and A. E. Raftery (2007).
Strictly proper scoring rules, prediction, and estimation.
\emph{Journal of the American Statistical Association}, 102(477), 359--378.
\href{https://doi.org/10.1198/016214506000001437}{doi:10.1198/016214506000001437}.
\bibitem{sportsdataverse}
SportsDataverse. \emph{sportsdataverse-data}, ESPN NBA schedules, team box scores,
and play-by-play releases, ending years 2022--2026.
\url{https://github.com/sportsdataverse/sportsdataverse-data}.
Exact URLs and hashes are recorded in collection manifests.
\bibitem{espn}
ESPN. NBA public game-summary responses.
\url{https://site.api.espn.com/apis/site/v2/sports/basketball/nba/summary}.
Event URLs and source hashes are recorded with primary replacements.
\bibitem{geonames}
GeoNames. \emph{cities15000 geographical database}.
\url{https://download.geonames.org/export/dump/}.
Creative Commons Attribution 4.0.
\bibitem{scipy}
SciPy. \emph{scipy.optimize.differential\_evolution}, version 1.18.1.
\url{https://docs.scipy.org/doc/scipy/reference/generated/scipy.optimize.differential_evolution.html}.
\bibitem{injuries}
NBA. Official injury-report publication policy and historical report PDFs.
\url{https://official.nba.com/nba-injury-report-2026-27-season/}.
Individual publication timestamps and hashes are retained with research sources.
\end{thebibliography}
\endgroup
\end{document}
